\documentclass[10pt,a4paper]{amsart}
\usepackage[margin=19mm]{geometry}
\usepackage{amsmath,amssymb,mathtools}
\usepackage[scr=rsfs,cal=euler]{mathalfa}
\usepackage{xfrac}
\usepackage[unicode,hidelinks,bookmarks=false]{hyperref}
\newcommand{\ie}{i.e.\ }
\newcommand{\cf}{cf.\ }
\newcommand{\resp}{respectively\ }
\newcommand{\Subsection}{\subsection}
\providecommand{\nobreakdash}{\nobreak\hbox{-}}
\setlength{\emergencystretch}{2em}
\multlinegap0pt
\usepackage{amssymb}
\usepackage{dsfont}
\usepackage{xcolor}
\newtheorem{lemma}{Lemma}
\newcommand{\R}{\mathbb{R}}
\newcommand{\sym}{\mathrm{sym}}
\title{Full flexibility of the Monge-Amp\`ere system in codimension $d_*-d+1$}
\author{Wentao Cao, Jonas Hirsch, Dominik Inauen, Marta Lewicka}
\date{Source: arXiv:2603.28909v1 (CC BY 4.0)}
\begin{document}
\maketitle

\noindent\emph{Source and licence.} Full flexibility of the Monge-Amp\`ere system in codimension $d_*-d+1$. Version arXiv:2603.28909v1 (CC BY 4.0), distributed by arXiv under the Creative Commons Attribution 4.0 International licence, http://creativecommons.org/licenses/by/4.0/, as recorded in the arXiv licence metadata for this exact version and shown on the abstract page https://arxiv.org/abs/2603.28909v1. Full licence notice: \texttt{licenses/arxiv-2603.28909-CC-BY-4.0.txt}.\par\smallskip

\noindent\textit{Editorial scope.} Counted unit: the article's lemma lem\_IBP (source0.tex lines 554--582). The article's complete original proof of it is delivered with it (source0.tex lines 583--627, one \texttt{proof} environment of 45 lines). No mathematics is written, repaired or re-derived: the statement and the proof are the article's own lines, in the article's own notation, and no retained line is substituted or re-flowed.  Retained uncounted as the source-local context the proof needs: lines 485--500.  Nothing was omitted from any retained span, so the package carries no omission record.  No retained line contains a citation, so the wrapper carries no bibliography and no bibliography entry.  No retained line contains a figure environment, an image-inclusion command or drawing code, so no image is delivered and the package declares no asset dependency.  Editorial formatting only, in addition: an AMS \texttt{amsart} preamble that restates the article's own declarations and macros used by the retained lines, namely the environments \texttt{lemma} on separate counters, together with the standard packages those lines need.  The retained environments print in this document as Lemma 1 in order of appearance, whereas the article prints the same items with the numbering of its own full document.  The complete native source of the article as distributed in the arXiv e-print for this exact version is bundled unchanged as \texttt{sources/197476/source0.tex}.
\begin{lemma}\label{lem_dec_def}
For all $i,j=1\ldots d$, define the unit vectors:
$$\eta_{ij}\doteq\frac{e_i+e_j}{|e_i+e_j|} \in\R^d.$$
Then $\{\eta_{ij}\otimes \eta_{ij}\}_{i,j=1\ldots d,\;i\leq j}$ is a basis of the
linear space $\R^{d\times d}_\sym$, so that:
$$H=\sum_{i\leq j}\bar a_{ij}(H) \eta_{ij}\otimes \eta_{ij}
\quad\mbox{ for all } H\in \R^{d\times d}_\sym,$$
where $\big\{\bar a_{ij}:\R^{d\times d}_\sym\to \R\big\}_{i,j=1\ldots d,\;i\leq j}$
are the linear coordinate functions. Moreover, denoting:
\begin{equation}\label{H_0}
H_0 \doteq \sum_{i\leq j}\eta_{ij}\otimes \eta_{ij},
\end{equation}
there exists $r_0\in (0, 1)$, depending only on $d$, such that:
$$|\bar a_{ij} (H)-1|\leq\frac{1}{2} \quad \mbox{ for all } H\in
B(H_0,r_0)\subset \R^{d\times d}_{\sym} \quad \mbox{and all }\; i \leq j.$$
\end{lemma}

\begin{lemma}\label{lem_IBP}
Given  $H\in\mathcal{C}^{k+1}(\R^d,\R^{d\times d}_\sym)$, $\lambda>0$,
$j=1\ldots d$ and $\Gamma_0\in\mathcal{C}(\R,\R)$, there holds:
\begin{equation}\label{ibp}
\begin{split}
\frac{\Gamma_0(\lambda t_{\eta_{1j}})}{\lambda}H = & \; (-1)^{k+1}\frac{\Gamma_{k+1}(\lambda
  t_{\eta_{1j}})}{\lambda^{k+2}} \sym\nabla L^{j}_k(H)  \\ & + \sym\nabla
\Big(\sum_{i=0}^k(-1)^i\frac{\Gamma_{i+1}(\lambda t_{\eta_{1j}})}{\lambda^{i+2}}L_i^j(H)\Big) 
+ \sum_{i=0}^{k}(-1)^i\frac{\Gamma_i(\lambda t_{\eta_{1j}})}{\lambda^{i+1}} P_i^j(H),
\end{split}
\end{equation}
where $t_{\eta_{ij}}=\langle x, \eta_{1j}\rangle$ with
$\eta_{1j} = \frac{e_1+e_j}{|e_1+e_j|}$. The functions $\Gamma_i\in
\mathcal{C}^{i}(\R,\R)$ satisfy the recursive definition: 
\begin{equation}\label{recu_ibp}
\Gamma_{i+1}' = \Gamma_{i}  \quad\mbox{ for all } \; i =0\ldots k,
\end{equation}
while $L_i^j(H)\in\mathcal{C}^{k+1-i}(\R^d,\R^d)$ and
$P_i^j(H)\in\mathcal{C}^{k+1-i}(\R^d,\R^{d\times d}_\sym)$ have the following properties:
\begin{itemize}
\item[(i)] the entries of $L_i^j(H)$ and $P_i^j(H)$ depend linearly on
  $H$ and consist of the linear combinations of the entries
  of $\nabla^{(i)}H$, 
\item[(ii)] $P_i^j(H)\in\mathrm{span}\big(\eta_{st}\otimes
  \eta_{st};\; 2\leq s\leq t\big)$, so that, recalling the
  decomposition in Lemma \ref{lem_dec_def}, there holds: $\bar
  a_{1t}(P_i^j(H))=0$ for all $t=1\ldots d$.
\end{itemize}
\end{lemma}

\begin{proof}
We assume that the matrix field $H$ is smooth and 
validate (\ref{ibp}) inductively for any $k\geq 0$, constructing the fields
$\{L_i^j =L_i^j(H)\}_{i\geq 0}$ and $\{P_i^j=P_i^j(H)\}_{i\geq 0}$ with the properties stated in (i), (ii). 
At $k=0$, we may check directly that these are valid in:
\begin{equation}\label{ibp0}
\begin{split}
& \frac{\Gamma_0(\lambda t_{\eta_{1j}})}{\lambda}H = \frac{\Gamma_0(\lambda t_{\eta_{1j}})}{\lambda}
\sym\big(L_0^{j}\otimes\eta_{ij}\big) + \frac{\Gamma_0(\lambda t_{\eta_{1j}})}{\lambda}P_0^j
\\ & \hspace{1.95cm} 
= -\frac{\Gamma_{1}(\lambda t_{\eta_{1j}})}{\lambda^{2}} \sym\nabla L^{j}_0  + \sym\nabla
\Big(\frac{\Gamma_{1}(\lambda t_{\eta_{1j}})}{\lambda^{2}}L_0^j\Big) 
+ \frac{\Gamma_0(\lambda t_{\eta_{1j}})}{\lambda} P_0^j, \\
& \mbox{where: } \;  L_0^j = 2\sqrt{2}He_1 - \sqrt{2}H_{11}(e_1+e_j),
\quad P_0^j = H-\sym(L_0^j\otimes \eta_{1j}).
\end{split}
\end{equation}
For the induction step, assume that $\{L_i^j, P_i^j\}_{i=0}^{k-1}$
are defined and (\ref{ibp}) holds at some $k-1\geq 0$:
\begin{equation}\label{ibpk-1}
\begin{split}
\frac{\Gamma_0(\lambda t_{\eta_{1j}})}{\lambda}H = & \; (-1)^{k}\frac{\Gamma_{k}(\lambda
  t_{\eta_{1j}})}{\lambda^{k+1}} \sym\nabla L^{j}_{k-1}  \\ & + \sym\nabla
\Big(\sum_{i=0}^{k-1}(-1)^i\frac{\Gamma_{i+1}(\lambda t_{\eta_{1j}})}{\lambda^{i+2}}L_i^j\Big) 
+ \sum_{i=0}^{k-1}(-1)^i\frac{\Gamma_i(\lambda t_{\eta_{1j}})}{\lambda^{i+1}} P_i^j.
\end{split}
\end{equation}
We now apply (\ref{ibp0}) to $H=\sym\nabla L_{k-1}^j$ and write:
\begin{equation*}
\begin{split}
& \frac{\Gamma_k(\lambda t_{\eta_{1j}})}{\lambda} \sym\nabla L_{k-1}^j
= -\frac{\Gamma_{k+1}(\lambda t_{\eta_{1j}})}{\lambda^{2}} \sym\nabla
L^{j}_k  + \sym\nabla
\Big(\frac{\Gamma_{k+1}(\lambda t_{\eta_{1j}})}{\lambda^{2}}L_k^j\Big) 
+ \frac{\Gamma_k(\lambda t_{\eta_{1j}})}{\lambda} P_k^j, \\
& \mbox{where: } \;  L_k^j = 2\sqrt{2}(\sym\nabla L_{k-1}^j)e_1 -
\sqrt{2}\langle \partial_1L_{k-1}^j, e_1\rangle (e_1+e_j), \\ & 
\hspace{1.35cm} P_k^j = \sym\nabla L_{k-1}^j - \sym(L_k^j\otimes \eta_{1j}).
\end{split}
\end{equation*}
Multiplying the above by $(-1)/\lambda^k$ and replacing with the
first term in the right hand side of (\ref{ibpk-1}), provides the
identity (\ref{ibp}). The recursive formulas for $L_k^j, P_k^j$ imply
the propagation of the properties (i), (ii). The proof is done. 
\end{proof}

\end{document}
